\documentclass[a4paper, serif, 10pt]{moderncv}

%% moderncv
% color and style
\moderncvcolor{black}
\moderncvstyle{banking}

%% page layout/margins
\usepackage[top=2.5cm, bottom=2.5cm, left=1.5cm, right=1.5cm]{geometry}

\nopagenumbers{}

%% Babel config for Norwegian language
% ref:
% - https://latex3.github.io/babel/guides/locale-norwegian.html
\usepackage[norsk]{babel}

%% fontspec
\usepackage{fontspec}

\IfFontExistsTF{Linux Libertine}{
  \setmainfont[Ligatures={TeX, Common}, Numbers=OldStyle]{Linux Libertine}
}{}
\IfFontExistsTF{Linux Libertine O}{
  \setmainfont[Ligatures={TeX, Common}, Numbers=OldStyle]{Linux Libertine O}
}{}

%% CTeX
% CJK support, used to show CJK characters in the resume
%
% - heading=false: disable ctex heading
% - punct=kaiming: use kaiming punctuations styles for CJK
% - scheme=plain: use plain scheme, do not override \normalsize font size
% - space=auto: space settings for CJK characters
%
% CJK fonts are resolved per locale language, e.g. Noto Serif/Sans CJK JP
% for Japanese resumes, with fallback to the next region if not installed
%
% ref:
% - http://ctan.mirrorcatalogs.com/language/chinese/ctex/ctex.pdf
\usepackage[UTF8, heading=false, punct=kaiming, scheme=plain, space=auto]{ctex}

\IfFontExistsTF{Noto Serif CJK SC}{
  \setCJKmainfont{Noto Serif CJK SC}
}{}
\IfFontExistsTF{Noto Sans CJK SC}{
  \setCJKsansfont{Noto Sans CJK SC}
}{}

%% line spacing
\usepackage{setspace}
\setstretch{1.125}

%% URL styling - use normal text instead of monospace
\urlstyle{same}

%% Patch \httplink and \httpslink to handle full URLs
% The original moderncv macros always prepend http:// or https:// to URLs.
% This causes issues when the URL already has a protocol (e.g., https://example.com)
% resulting in malformed URLs like https://https://example.com.
% This patch checks if the URL already contains "://" and uses it directly if so.
% Note: We use \str_set:Nx (x-expansion) to expand macros like \@homepage first.
\ExplSyntaxOn
\str_new:N \g_yamlresume_url_str

\RenewDocumentCommand{\httpslink}{O{}m}{
  \str_set:Nx \g_yamlresume_url_str {#2}
  \str_if_in:NnTF \g_yamlresume_url_str {://}
    { \url{#2} }
    { \url{https://#2} }
}

\RenewDocumentCommand{\httplink}{O{}m}{
  \str_set:Nx \g_yamlresume_url_str {#2}
  \str_if_in:NnTF \g_yamlresume_url_str {://}
    { \url{#2} }
    { \url{http://#2} }
}
\ExplSyntaxOff

\name{Kristoffer A. Løken}{}
\title{Senior finansanalytiker med fokus på verdsettelse, risiko og kapitalmarkeder}
\phone[mobile]{+47 912 34 567}
\email{kristoffer.loken@fjordkapital.no}
\homepage{https://www.kristofferloken.no}

\address{Storgata 45, Oslo, Oslo, Norge, 0182}{}{}

\extrainfo{{\small \faLinkedin}\ \href{https://www.linkedin.com/in/kristoffer-loken}{@kristoffer-loken} {} {} {} • {} {} {} 
{\small \faGithub}\ \href{https://github.com/kloken-finance}{@kloken-finance}}

\begin{document}

\maketitle

\section{Grunnleggende informasjon}

\cvline{}{\begin{itemize}
\item Finansanalytiker med mastergrad i finansiell økonomi og flere års erfaring fra bank og kapitalforvaltning
\item Sterk kompetanse i verdsettelse, kredittanalyse, makroøkonomi og risikostyring
\item Erfaren med å bygge modeller og automatiserte rapporter i Excel, SQL, Python og Power BI
\item Opptatt av å omsette analyse til tydelige anbefalinger for beslutningstakere
\end{itemize}}
\section{Utdanning}

\cventry{aug. 2015–juni 2017}
        {Master, Finansiell økonomi, Poeng: 3,9/4,0}
        {Norges Handelshøyskole}
        {\url{https://www.nhh.no}}
        {}
        {\begin{itemize}
\item Fordypning i verdsettelse, kapitalmarkeder og kvantitativ analyse
\item Masteroppgave om kredittrisiko i nordiske obligasjonsfond
\item Gjennomsnittskarakter 3,9 av 4,0
\end{itemize}
\textbf{Kurs}: Selskapsfinans, Verdsettelse og investeringsanalyse, Risikostyring og derivater, Regnskapsanalyse, Makroøkonomi, Porteføljeteori, Kvantitative metoder, ESG og bærekraftig finans}

\cventry{aug. 2012–juni 2015}
        {Bachelor, Samfunnsøkonomi, Poeng: 3,7/4,0}
        {Universitetet i Oslo}
        {\url{https://www.uio.no}}
        {}
        {\begin{itemize}
\item Grunnleggende fordypning i samfunnsøkonomi og statistikk
\item Prosjektoppgave om renteutvikling og boligmarkedet i Norge
\end{itemize}
\textbf{Kurs}: Mikroøkonomi, Makroøkonomi, Statistikk, Økonometri, Finansmarkeder, Regnskap og budsjettering}
\section{Arbeidserfaring}

\cventry{jan. 2021–Nå}
        {Senior finansanalytiker}
        {Fjordkapital Forvaltning}
        {\url{https://www.fjordkapital.no}}
        {}
        {\begin{itemize}
\item Ansvarlig for verdsettelsesmodeller (DCF, DDM) for nordiske aksjer og obligasjoner
\item Utarbeider kvartalsvise investeringsanbefalinger til porteføljeforvaltere og risikokomité
\item Bygger automatiserte rapporter i Power BI og Python som reduserer manuell rapportering med 30 \%
\item Overvåker makro-, rente- og kredittmarkeder og bidrar i ESG-analyse
\item Mentorerer to junioranalytikere og kvalitetssikrer modeller og investeringsnotater
\end{itemize}
\textbf{Nøkkelord}: Verdsettelse, DCF, ESG, Porteføljestyring, Power BI, Risikoanalyse}

\cventry{sep. 2017–des. 2020}
        {Finansanalytiker}
        {Nordlys Bank}
        {\url{https://www.nordlysbank.no}}
        {}
        {\begin{itemize}
\item Analyserte kredittporteføljer og laget månedlige risikorapporter til ledelsen
\item Utviklet stresscenarier og makroøkonomiske prognoser for kapitalplanlegging
\item Forbedret datakvalitet i SQL-baserte rapporteringsløsninger
\item Presenterte funn for kredittkomité og bidro til beslutningsgrunnlag
\end{itemize}
\textbf{Nøkkelord}: Kredittanalyse, Stresscenarier, SQL, Rapportutvikling, Kapitalplanlegging}

\cventry{juni 2016–aug. 2016}
        {Sommeranalytiker}
        {Oslo Analytics Partners}
        {\url{https://www.osloanalytics.no}}
        {}
        {\begin{itemize}
\item Bygget markedspaneler og scenarioanalyser for energisektoren
\item Støttet senioranalytikere med datainnsamling og kvalitetssikring
\item Laget Excel-modeller for kostnads- og inntektsprognoser
\end{itemize}
\textbf{Nøkkelord}: Excel, Markedsanalyse, Energi, Modellering}
\section{Språk}

\cvline{Norsk}{Morsmål eller tospråklig nivå \hfill \textbf{Nøkkelord}: Morsmål}
\cvline{Engelsk}{Fullt profesjonelt nivå \hfill \textbf{Nøkkelord}: C1, Forretningsengelsk}
\cvline{Tysk}{Begrenset arbeidskunnskap \hfill \textbf{Nøkkelord}: A2, Grunnleggende}
\section{Ferdigheter}

\cvline{Finansiell modellering}{Ekspert \hfill \textbf{Nøkkelord}: DCF, DDM, LBO, Scenarioanalyse, Excel, VBA}
\cvline{Dataanalyse og visualisering}{Avansert \hfill \textbf{Nøkkelord}: SQL, Python, Pandas, Power BI, Tableau, R}
\cvline{Risiko og kapitalstyring}{Avansert \hfill \textbf{Nøkkelord}: Kredittrisiko, Markedsrisiko, Stresscenarier, Basel III, ICAAP}
\cvline{Investeringsanalyse}{Avansert \hfill \textbf{Nøkkelord}: Aksjeanalyse, Obligasjoner, Makroøkonomi, Porteføljeteori, ESG}
\cvline{Kommunikasjon og rådgivning}{Ekspert \hfill \textbf{Nøkkelord}: Investeringsnotater, Presentasjoner, Interessentstyring, Mentoring}
\section{Utmerkelser}

\cventry{des. 2023}
        {Fjordkapital Forvaltning}
        {Årets analytiker}
        {}
        {}
        {Tildelt for fremragende analysearbeid og bidrag til forbedret investeringsprosess i nordiske aksjer.}

\cventry{juni 2017}
        {Norges Handelshøyskole}
        {Beste masteroppgave i finansiell økonomi}
        {}
        {}
        {Pris for masteroppgave om kredittrisiko og likviditet i nordiske obligasjonsfond.}
\section{Sertifikater}

\cventry{aug. 2020}
        {CFA Institute}
        {CFA nivå II}
        {\url{https://www.cfainstitute.org/programs/cfa}}
        {}
        {}

\cventry{mars 2019}
        {Corporate Finance Institute}
        {FMVA – finansiell modellering og verdsettelse}
        {\url{https://corporatefinanceinstitute.com/certifications/fmva/}}
        {}
        {}
\section{Publikasjoner}

\cventry{nov. 2018}
        {Kredittrisiko i nordiske obligasjonsfond}
        {Nordisk tidsskrift for finans}
        {\url{https://www.nordiskfinans.no/artikler/kredittrisiko-i-nordiske-obligasjonsfond}}
        {}
        {\begin{itemize}
\item Analyserer hvordan likviditet og kredittspreader påvirker avkastning i nordiske obligasjonsfond
\item Funn tyder på at stresstester bør inkludere likviditetspremier i større grad
\end{itemize}}
\section{Prosjekter}

\cventry{jan. 2022–okt. 2023}
        {Modell for å integrere ESG-indikatorer i verdsettelse og risikovurdering}
        {ESG-risikomodell for nordiske aksjer}
        {\url{https://www.fjordkapital.no/innsikt/esg-risikomodell}}
        {}
        {\begin{itemize}
\item Utviklet scoringsrammeverk for klimarisiko, styring og sosiale forhold
\item Integrerte modellen i kvartalsvis investeringsprosess
\item Presenterte resultater for investeringskomité og kunder
\end{itemize}
\textbf{Nøkkelord}: ESG, Klimarisiko, Verdsettelse, Power BI}

\cventry{feb. 2019–nov. 2020}
        {Python-basert rapportløsning som samler kredittdata og nøkkeltall}
        {Automatisert månedsrapport for kredittportefølje}
        {\url{https://www.nordlysbank.no/om-oss/rapportering}}
        {}
        {\begin{itemize}
\item Reduserte manuell rapporttid fra tre dager til fire timer
\item Sikret konsistent datakvalitet og sporbarhet i rapporteringen
\item Bygget dashboards for ledelsen i Power BI
\end{itemize}
\textbf{Nøkkelord}: Python, SQL, Power BI, Automatisering}
\section{Interesser}

\cvline{Friluftsliv}{Fjellturer, Langrenn, Kajakk}
\cvline{Samfunn og økonomi}{Makroøkonomi, Bærekraft, Litteratur}
\section{Frivillig arbeid}

\cventry{jan. 2022–des. 2024}
        {Mentor og fagansvarlig}
        {Norsk Forum for Finansanalyse}
        {\url{https://www.norskfinansanalyse.no}}
        {}
        {\begin{itemize}
\item Mentorerer studenter og nyutdannede innen finansiell analyse
\item Bidrar til fagkvelder om verdsettelse, makro og ESG
\item Koordinerer case-konkurranse for studenter
\end{itemize}}
    
\end{document}